Metarhia Protocol являє собою протокол взаємодії мережевих додатків і їх
компонентів, що надає можливості для двостороннього асинхронного обміну даними,
мультиплексування багатьох паралельних каналів стрімінгу даних,
\gls{RPC}-викликів та потоків подій, на які можна підписуватися, і може прозоро
для прикладного коду обробляти короткочасні обриви з'єднання, повністю
відновлюючи сесію. Також він реалізує механізм аутентифікації з'єднань,
опціональне стиснення даних і підтримує різні формати серіалізації даних.

\subsubsection{Визначення}

Розглянемо назви сутностей та абстракцій, що використовуються в даній специфікації
та програмному коді, що її реалізує.

\emph{Transport} (<<транспорт>>) - транспортний протокол або інший механізм
коммунікації, що надає повнодуплексне двостороннє з'єднання з гарантіями
доставки та збереження порядку відправлення даних на приймаючій стороні, і
якому Metarhia Protocol делегує низькорівневу передачу потоків бінарних даних.

\emph{Transport connection} (<<транспортне з'єднання>>) - сокет, з'єднання,
або інший специфічний для транспорту канал комунікації.

\emph{Metarhia Protocol connection} (або просто \emph{connection} без уточнень
- <<з'єднання>>) - абстракція над з'єднанням транспортного протоколу, що
управляється реалізацією Metarhia Protocol, яка приховує деталі реалізації
транспорту та надає функціональність прикладного протоколу користувачеві (тобто
розробнику прикладних додатків). Одному з'єднанню Metarhia Protocol відповідає
одне основне транспортне з'єднання (що використовується для \gls{RPC}), проте
реалізація Metarhia Protocol може за необхідності внутрішньо відкривати
додаткові з'єднання.

\emph{Chunk} (буквально <<шматок>>) - одиниця даних протоколу Metarhia
Protocol, що складається з заголовків і опціонального навантаження (payload).
Всі дані, що передаються протоколом, розбиваються на чанки.

\emph{Channel} (<<канал>>) - це множина чанків, що передаються через певне
з'єднання та ідентифікуються серед інших певним числовим значенням, що є
унікальним серед активних каналів цього з'єднання. Канали надають
функціональність мультиплексування паралельних потоків даних.

\emph{Message} (<<повідомлення>>) - це канал, що характеризується коротким
часом існування і невеликою та обмеженою кількістю чанків (більшість
повідомлень не перевищують за обсягом один чанк), що буферизуються в пам'яті до
повного отримання повідомлення, після чого обробляються як єдине ціле. До
повідомлень відносяться \gls{RPC}-виклики та віддалені події.

\emph{Stream} (<<потік>>) - канал, що характеризується довільним часом
існування (обмеженим лише часом існування самого з'єднання) та невизначеною
кількістю чанків, дані яких обробляються прикладним кодом одразу по мірі
доступності.

\emph{Session} (<<сесія>>) - стійка асоціація між додатками на обох кінцях
з'єднання. Сесії можуть бути анонімними та аутентифікованими.  Канали є
прив'язаними до відповідних сесій.

\emph{Tunnel} (<<тунель>>) - набір сесій і каналів, що належать з'єднанню.
Ця абстракція реалізує прозоре перепідключення в разі втрати транспортного
з'єднання.

Ключові слова \emph{<<повинні>>}, \emph{<<слід>>}, \emph{<<можуть>>} і
\emph{<<опціонально>>} в цьому підрозділі відповідають англомовним аналогам
<<MUST>>, <<SHOULD>>, <<MAY>> і <<OPTIONAL>> в тому сенсі, в якому вони
визначені в RFC 2119~\cite{rfc2119}.

\subsubsection{Підтримка транспортних протоколів}

Реалізації Metarhia Protocol, що в основному націлені на використання в
серверному оточенні, \emph{повинні} підтримувати протоколи \gls{TCP}
(див.~\cite{rfc793}), \gls{TLS} (див.~\cite{rfc5246}), WebSocket
(див.~\cite{rfc6455}) і WebSocket over \gls{TLS} в якості транспортних.
Реалізації Metarhia Protocol, що призначені для використання в серверному
оточенні, \emph{можуть} також підтримувати додаткові транспорти, наприклад,
Unix domain сокети для міжпроцесної взаємодії в рамках фізичної машини за
допомогою тих самих абстракцій (що особливо актуально для додатків з
мікросервісною архітектурою).

Реалізаціям Metarhia Protocol, що розроблені спеціально для використання
клієнтськими додатками, за можливості \emph{слід} підтримувати транспортні
протоколи \gls{TCP} і \gls{TLS}.  В тих випадках, коли це неможливо (наприклад,
реалізації для веб-додатків, що виконуються в браузері), ці реалізації
\emph{повинні} підтримувати транспортні протоколи WebSocket і WebSocket over
\gls{TLS}, в іншому випадку підтримка WebSocket є \emph{опціональною}.  Іншими
словами, принаймні одна з пар <<\gls{TCP} і \gls{TLS}>> або <<WebSocket і
WebSocket over \gls{TLS}>> \emph{повинна} підтримуватися, з перевагою в бік
\gls{TCP} і \gls{TLS}.  Клієнтські реалізації також \emph{можуть} підтримувати
інші транспортні протоколи, якщо в цьому буде практичний сенс.

\subsubsection{Стани з'єднання}

Діаграма переходів між станами наведена в документі ІАЛЦ.467100.006~Д2.

\paragraph{AWAITING\_HANDSHAKE}

Транспортне з'єднання було відкрито, проте рукостискання (handshake) ще не відбулося
і не була встановлена сесія.

Клієнт надсилає чанк \emph{Protocol Handshake}. Коли рукостискання успішно
завершується, з'єднання переходить в стан \emph{AWAITING\_TUNNEL}.

Сервер очікує чанк \emph{Protocol Handshake} від клієнта. Коли рукостискання
успішно завершується, з'єднання переходить в стан \emph{AWAITING\_TUNNEL}.

\paragraph{AWAITING\_TUNNEL}

Клієнт надсилає чанк \emph{Open Tunnel}. Сервер відповідає або чанком \emph{New
Tunnel}, або, в разі відновлення з'єднання, чанком \emph{State
Synchronization}, на який клієнт відповідає також чанком \emph{State
Synchronization}, і обидві сторони обмінюються буферизованими даними, що не
передавалися за період відсутності з'єднання. Після цього обидва з'єднання
(тобто на стороні клієтського та серверного додатків) переходять у стан
\emph{NORMAL}.

\paragraph{NORMAL}

Це основний режим роботи протоколу. Вся комунікація відбувається за рахунок
каналів і чанків ping/pong. У випадку мережевої помилки з'єднання переходять у
стан \emph{NETWORK\_CONN\_LOST}.

\paragraph{NETWORK\_CONN\_LOST}

Клієнт буферизує вихідні чанки та намагається перепідключитися до сервера.
Якщо це вдається, то з'єднання переходить у стан \emph{AWAITING\_HANDSHAKE}.

Сервер буферизує вихідні чанки та очікує нове з'єднання від клієнта.

\subsubsection{Різновиди одиниць даних типу Chunk}

Кожен чанк, що передається в стані з'єднання \emph{NORMAL}, починається з поля
довжиною один октет, що вказує на тип чанка.  Це значення \emph{повинно} бути
одним із наведених у таблиці~\ref{tbl:chunk-types}.

\begin{table}[htbp]
  \captionbox{\label{tbl:chunk-types}}{%
    \begin{tabu} to \textwidth { | X[l] | X[l] | }
      \hline
      Назва типу & Значення першого октета \\
      \hline
      PING & 0 \\
      \hline
      PONG & 1 \\
      \hline
      MESSAGE\_PREAMBLE & 2 \\
      \hline
      STREAM\_PREAMBLE & 3 \\
      \hline
      DATA\_CHUNK & 4 \\
      \hline
    \end{tabu}
  }
\end{table}

\subsubsection{Формати одиниць даних типу Chunk}

Зауваження: Metarhia Protocol використовує порядок байтів little-endian.

\paragraph{Protocol Handshake}

У таблиці~\ref{tbl:protocol-handshake} подано формат чанку <<Protocol Handshake>>.

\begin{table}[htbp]
  \captionbox{\label{tbl:protocol-handshake}}{%
    \begin{tabu} to \textwidth { | X[l] | X[l] | }
      \hline
      Поле & Кількість бітів \\
      \hline
      Version & 16 \\
      \hline
      Encryption & 16 \\
      \hline
      \multicolumn{2}{|l|}{Payload} \\
      \hline
    \end{tabu}
  }
\end{table}

Поле Version вказує на версію протоколу (тобто, 1). Єдиним можливим значенням
поля Encryption на даний момент є 0, а Payload відсутній.

Якщо в подальшому будуть додані інші значення поля Encryption, то вони також
можуть вказувати на необхідність вибору сторонами складнішого процесу
рукостискання (наприклад, для обміну ключами).

\paragraph{Open Tunnel}

У таблиці~\ref{tbl:protocol-opentunnel} подано формат чанку <<Open Tunnel>>.

\begin{table}[htbp]
  \captionbox{\label{tbl:protocol-opentunnel}}{%
    \begin{tabu} to \textwidth { | X[l] | X[l] | }
      \hline
      Поле & Кількість бітів \\
      \hline
      Token & 256 \\
      \hline
    \end{tabu}
  }
\end{table}

Token - ідентифікатор тунелю і його сесійний ключ, що служить для повторної
аутентифікації при відновленні сесії.  Значення 0 має особливе значення і
служить для створення нового тунелю замість існуючого.

\paragraph{New Tunnel}

У таблиці~\ref{tbl:protocol-newtunnel} подано формат чанку <<New Tunnel>>.

\begin{table}[htbp]
  \captionbox{\label{tbl:protocol-newtunnel}}{%
    \begin{tabu} to \textwidth { | X[l] | X[l] | }
      \hline
      Поле & Кількість бітів \\
      \hline
      Token & 256 \\
      \hline
    \end{tabu}
  }
\end{table}

Token - це отриманий з криптографічно стійкого джерела 32-байтний випадковий
рядок, що служить ідентифікатором тунелю і його секретним ключем. Token не може
бути рівним нулю.

\paragraph{State Synchronization}

У таблиці~\ref{tbl:protocol-statesync} подано формат чанку <<State Synchronization>>.

\begin{table}[htbp]
  \captionbox{\label{tbl:protocol-statesync}}{%
    \begin{tabu} to \textwidth { | X[l] | X[l] | }
      \hline
      Поле & Кількість бітів \\
      \hline
      LastPingId & 32 \\
      \hline
      ChunksCount & 32 \\
      \hline
    \end{tabu}
  }
\end{table}

LastPingId - це ідентифікатор останнього Ping-чанку, що отримала сторона, що
надіслала даний чанк State Synchronization. ChunksCount - кількість чанків,
що вона отримала з того моменту.

\paragraph{Ping}

У таблиці~\ref{tbl:protocol-ping} подано формат чанку <<Ping>>.

\begin{table}[htbp]
  \captionbox{\label{tbl:protocol-ping}}{%
    \begin{tabu} to \textwidth { | X[l] | X[l] | }
      \hline
      Поле & Кількість бітів \\
      \hline
      ChunkType & 8 \\
      \hline
      PingId & 32 \\
      \hline
    \end{tabu}
  }
\end{table}

Значення ChunkType у Ping-чанків дорівнює константі PING (див.
таблицю~\ref{tbl:chunk-types}). PingId - ціле число, що монотонно зростає.

\paragraph{Pong}

У таблиці~\ref{tbl:protocol-pong} подано формат чанку <<Pong>>.

\begin{table}[htbp]
  \captionbox{\label{tbl:protocol-pong}}{%
    \begin{tabu} to \textwidth { | X[l] | X[l] | }
      \hline
      Поле & Кількість бітів \\
      \hline
      ChunkType & 8 \\
      \hline
      PingId & 32 \\
      \hline
    \end{tabu}
  }
\end{table}

Значення ChunkType у Pong-чанків дорівнює константі PONG (див.
таблицю~\ref{tbl:chunk-types}). PingId - ідентифікатор відповідного чанку
Ping.

\paragraph{Channel Preamble}

У таблиці~\ref{tbl:protocol-chan-preamble} подано загальну частину структури,
що є першим чанком всіх каналів. Структури, що представляють перші чанки
конкретних типів каналів - Message Preamble
(таблиця~\ref{tbl:protocol-msg-preamble}) і Stream Preamble, - наслідують і
розширяють її.

Id і Compression - загальні поля, що характерні для обох видів каналів.
Структура Stream Preamble не містить додаткових полів, тому
таблиця~\ref{tbl:protocol-chan-preamble}, фактично, описує її також.  Додаткові
поля структури Message Preamble розміщуються на місці поля
MessagePreambleReserved абстрактної структури Channel Preamble.

Поле ChunkType структури Message Preamble дорівнює константі MESSAGE\_PREAMBLE,
а поле ChunkType структури Stream Preamble дорівнює константі STREAM\_PREAMBLE
(див. таблицю~\ref{tbl:chunk-types}).

\begin{table}[htbp]
  \captionbox{\label{tbl:protocol-chan-preamble}}{%
    \begin{tabu} to \textwidth { | X[l] | X[l] | }
      \hline
      Поле & Кількість бітів \\
      \hline
      ChunkType & 8 \\
      \hline
      Id & 32 \\
      \hline
      Compression & 8 \\
      \hline
      MessagePreambleReserved & 16 \\
      \hline
      SessionId & 64 \\
      \hline
    \end{tabu}
  }
\end{table}

Поле Id - це ідентифікатор каналу в з'єднанні. Щоб уникнути необхідності
синхронізувати лічильники на двох сторонах з'єднання і можливих колізій,
старший біт поля Id маскується і завжди дорівнює 0 для каналів, створених
клієнтом, і дорівнює 1 для каналів, створених сервером. Іншими словами, можна
сказати, що дане поле є знаковим цілим числом, і коректними його значеннями для
каналів, створених клієнтом, є $[0, 2^{31} - 1]$, а для каналів, створених
сервером, - $[{-2^{31}}, -1]$.

Значення поля Id \emph{повинно} бути унікальним серед активних у момент часу
каналів.

Поле Compression вказує, чи є дані наступних чанків з даними цього каналу
стисненими. Поле \emph{повинно} бути рівним одному із значень, поданих у
таблиці~\ref{tbl:protocol-compression}.

\begin{table}[htbp]
  \captionbox{\label{tbl:protocol-compression}}{%
    \begin{tabu} to \textwidth { | X[l] | X[l] | }
      \hline
      Значення поля Compression & Ефект \\
      \hline
      0 & Стиснення не використовується \\
      \hline
      1 & Дані стиснені алгоритмом GZip \\
      \hline
    \end{tabu}
  }
\end{table}

SessionId - це ідентифікатор сесії, в якій має бути відкритий канал.  Клієнт
отримує його у відповідь на успішне рукостискання з додатком, тобто на
повідомлення HandshakeRequest, або, іншими словами, авторизацію в ньому чи
створення анонімної сесії.

\paragraph{Message Preamble}

У таблиці~\ref{tbl:protocol-msg-preamble} подано формат чанку <<Message Preamble>>.
Див. також таблицю~\ref{tbl:protocol-chan-preamble}.

\begin{table}[htbp]
  \captionbox{\label{tbl:protocol-msg-preamble}}{%
    \begin{tabu} to \textwidth { | X[l] | X[l] | }
      \hline
      Поле & Кількість бітів \\
      \hline
      ChunkType & 8 \\
      \hline
      Id & 32 \\
      \hline
      Compression & 8 \\
      \hline
      Encoding & 8 \\
      \hline
      MessageType & 8 \\
      \hline
      SessionId & 64 \\
      \hline
    \end{tabu}
  }
\end{table}

Ця структура розширює загальну Channel Preamble і додає на зарезервоване
місце два нових поля: Encoding і MessageType.

Поле Encoding вказує формат серіалізації складних структур даних (наприклад,
аргументів \gls{RPC}-викликів у тілі повідомлень типу Call). Воно
\emph{повинно} бути рівним одному із значень, вказаних у
таблиці~\ref{tbl:protocol-encoding}.

Поле MessageType вказує на тип повідомлення і має бути одним із значень,
наведених у таблиці~\ref{tbl:protocol-msg-type}.

\begin{table}[htbp]
  \captionbox{\label{tbl:protocol-encoding}}{%
    \begin{tabu} to \textwidth { | X[l] | X[l] | }
      \hline
      Значення поля Encoding & Формат серіалізації \\
      \hline
      0 & JSTP \\
      \hline
      1 & JSON \\
      \hline
    \end{tabu}
  }
\end{table}

\begin{table}[htbp]
  \captionbox{\label{tbl:protocol-msg-type}}{%
    \begin{tabu} to \textwidth { | X[l] | X[l] | }
      \hline
      Значення поля MessageType & Тип повідомлення \\
      \hline
      0 & HandshakeRequest \\
      \hline
      1 & HandshakeResponse \\
      \hline
      2 & Event \\
      \hline
      3 & Call \\
      \hline
      4 & Callback \\
      \hline
      5 & Inspect \\
      \hline
      6 & InspectCallback \\
      \hline
    \end{tabu}
  }
\end{table}

\paragraph{Data Chunk}

У таблиці~\ref{tbl:protocol-data-chunk} подано формат чанку <<Data Chunk>>.

\begin{table}[htbp]
  \captionbox{\label{tbl:protocol-data-chunk}}{%
    \begin{tabu} to \textwidth { | X[l] | X[l] | }
      \hline
      Поле & Кількість бітів \\
      \hline
      ChunkType & 8 \\
      \hline
      ChannelId & 32 \\
      \hline
      Length & 16 \\
      \hline
      Flags & 8 \\
      \hline
    \end{tabu}
  }
\end{table}

Поле ChunkType цієї структури дорівнює константі DATA\_CHUNK
(див. таблицю~\ref{tbl:chunk-types}).

Поле ChannelId указує на те, до якого каналу належить цей чанк з даними.

Поле Length містить розмір наступного за цими заголовками блоку даних у байтах.

Поле Flags має структуру, що описана в таблиці~\ref{tbl:data-chunk-flags}.

\begin{table}[htbp]
  \captionbox{\label{tbl:data-chunk-flags}}{%
    \begin{tabu} to \textwidth { | X[l] | X[l] | }
      \hline
      Розряд & Значення булевого флагу \\
      \hline
      Біти 7-1 & Зарезервовано \\
      \hline
      Біт 0 & More \\
      \hline
    \end{tabu}
  }
\end{table}

Флаг More вказує, чи є в даному каналі наступний Data Chunk. Якщо значення біта
виставлено в 0, то даний Data Chunk є останнім і закриває собою канал.

Зарезервовані флаги \emph{повинні} бути виставлені в 0.

\subsubsection{Альтернативний метод швидкого сповіщення про події в системі}

Планується також використовувати UDP як допоміжний канал зв'язку в тих
випадках, коли порядок і гарантія доставки подій не важлива, прикріплені дані
малі за обсягом, а частота велика, тому передавати повідомлення типу Event
стандартним механізмом неефективно.

Для захисту основного каналу передачі даних є TLS, а конфіденційність і
цілісніть UDP-датаграм необхідно забезпечити на рівні прикладного протоколу.

Для цього будемо використовувати симетричне шифрування сесійним ключем
<<тунелю>> (див. таблицю~\ref{tbl:protocol-newtunnel}) з використанням
криптографічного примітиву \gls{AEAD}, що забезпечує одночасно і шифрування, і
криптографічний підпис, що засвідчує цілісність даних, а також дозволяє додати
до так званого \gls{AEAD}-контейнера невелику кількість незашифрованих (але
підписаних) даних, що в даному сценарії використання знадобиться для того, щоб
розрізняти, якому зі з'єднань (і, відповідно, <<тунелів>>) слід маршрутизувати
датаграму, якщо всі вони (або декілька з них) будуть використовувати один порт.

В якості конкретної реалізації концепції AEAD візьмемо сучасний алгоритм, що
базується на шифрі ChaCha20~\cite{chacha20} і алгоритмі аутентифікації
повідомлень (\gls{MAC}) Poly1305~\cite{poly1305}, з модифікаціями IETF, що
прийнято в якості стандарту RFC~7539~\cite{rfc7539}

І ChaCha20, і Poly1305 - це сучасні, надійні і дуже швидкі алгоритми, що
були детально досліджені багатьма криптографами, і зараз стрімко набирають
популярність на заміну більш класичним алгоритмам. В якості прикладів:

\begin{itemize}
\item
  Google використовував свою імплементацію цих алгоритмів для шифрування
  TLS-трафіку між своїми серверами і пристроями під управлінням Android, а
  також браузером Google Chrome, з 2014 року;
\item
  чернетка офіційного стандарту TLS~1.3 містить набір шифрів
  TLS\_CHACHA20\_POLY1305\_SHA256 і рекомендує його до використання~\cite{tls13}.
\end{itemize}

Версії алгоритмів ChaCha20, ChaCha20-Poly1305 і ChaCha20-Poly1305 \gls{AEAD}
від IETF, що описані в RFC~7539, відрізняються від оригінальних тим, що довжина
nonce-значення збільшена з 64 біт до 96 біт, через що розмір лічильника для
шифрування блоків зменшився з 64 біт до 32 біт, обмежуючи розмір одного
повідомлення 256 Гб замість $2^{64}$ байтів, але дозволяючи зашифрувати одним і
тим самим ключем значно більше різних повідомлень (як і в інших подібних
алгоритмах, не можна використовувати одну й ту ж пару ключа і нонса на різних
даних, інакше, перехопивши їх, атакуючий зможе провести криптоаналіз).

Таким чином, з кожним екземпляром структури Tunnel буде асоційовано два числа:
сесійний ключ довжиною 32 байти, що генерує сервер, і випадкове число
nonce довжиною 12 байтів, яке кожна сторона генерує для себе сама, після чого
маскує молодший біт: клієнт виставляє його в 0, а сервер - в 1.

Алгоритм використання \gls{AEAD} для шифрування датаграм показано на
рис.~\ref{fig:udp-encrypt-algo} і рис.~\ref{fig:udp-encrypt-flowchart}.
Алгоритм використання \gls{AEAD} для дешифрування датаграм показано на
рис.~\ref{fig:udp-decrypt-algo}, а повний алгоритм їх обробки після отримання
з зовнішньої мережі - в документі ІАЛЦ.467100.003~Д3.

\begin{figure}[!htbp]
  \begin{algorithm}[H]
    \SetKwData{Secret}{secret}
    \SetKwData{Tunnel}{tunnel}
    \SetKwData{Nonce}{nonce}
    \SetKwData{Message}{message}
    \SetKwData{Ciphertext}{container}
    \SetKwFunction{GetSecret}{GetSecret}
    \SetKwFunction{GetNonce}{GetNonce}
    \SetKwFunction{SetNonce}{SetNonce}
    \SetKwFunction{Encrypt}{AEAD\_ChaCha20\_Poly1305\_IETF\_Encrypt}
    \SetKwInOut{Input}{input}
    \SetKwInOut{Output}{output}

    \Input{%
      \Message\xspace - повідомлення для захищеної передачі в контексті
      мультиплексованого з'єднання \Tunnel; \Message~$=$~$(H, P)$, де $H$ -
      відкриті заголовки, $P$ - секретна частина}
    \Output{\Ciphertext\xspace - зашифроване повідомлення в \gls{AEAD}-контейнері}
    \BlankLine

    \Begin{
      \Secret$\leftarrow$ \GetSecret{\Tunnel}\;
      \Nonce$\leftarrow$ \GetNonce{\Tunnel}\;
      \SetNonce{\Tunnel, \Nonce$+$ $2$}\;
      \Ciphertext$\leftarrow$ \Encrypt{\Message, \Secret, \Nonce}\;
    }
  \end{algorithm}

  \caption{Алгоритм створення захищеного повідомлення у вигляді псевдокоду}%
  \label{fig:udp-encrypt-algo}
\end{figure}

\begin{figure}[!hbtp]
  \begin{algorithm}[H]
    \SetKwData{Secret}{secret}
    \SetKwData{Tunnel}{tunnel}
    \SetKwData{Message}{message}
    \SetKwData{Ciphertext}{container}
    \SetKwFunction{GetSecret}{GetSecret}
    \SetKwFunction{Decrypt}{AEAD\_ChaCha20\_Poly1305\_IETF\_Decrypt}
    \SetKwInOut{Input}{input}
    \SetKwInOut{Output}{output}

    \Input{\Ciphertext\xspace - повідомлення в \gls{AEAD}-контейнері,
      ключ для розшифрування якого є унікальним для мультиплексованого
      з'єднання \Tunnel; \Ciphertext~$=$~$(H, C, S)$, де $H$ - відкриті
      заголовки, $C$ - шифротекст секретної частини повідомлення, $S$ -
      \gls{MAC}-підпис}
    \Output{\Message\xspace - вихідне повідомлення}
    \BlankLine

    \Begin{
      \Secret$\leftarrow$ \GetSecret{\Tunnel}\;
      \Message$\leftarrow$ \Decrypt{\Ciphertext, \Secret}\;
    }
  \end{algorithm}

  \caption{Алгоритм відтворення захищеного повідомлення у вигляді псевдокоду}%
  \label{fig:udp-decrypt-algo}
\end{figure}

\begin{figure}[tbp]
  \centering

  \begin{tikzpicture}[node distance=3cm, font=\gostfont\small]
    \node (start) [flow:startstop] {Початок};
    \node (input) [flow:io, below of=start] {%
      Повідомлення, секретний ключ і останнє значення nonce};
    \node (encrypt) [flow:process, below of=input] {%
      Шифрування повідомлення в \gls{AEAD}-режимі};
    \node (upd-nonce) [flow:process, below of=encrypt] {%
      Збільшити значення nonce на 2};
    \node (output) [flow:io, below of=upd-nonce] {%
      Зашифроване повідомлення, нове значення nonce};
    \node (stop) [flow:startstop, below of=output] {Кінець};

    \draw [flow:arrow] (start) -- (input);
    \draw [flow:arrow] (input) -- (encrypt);
    \draw [flow:arrow] (encrypt) -- (upd-nonce);
    \draw [flow:arrow] (upd-nonce) -- (output);
    \draw [flow:arrow] (output) -- (stop);
  \end{tikzpicture}

  \caption{Блок-схема алгоритму створення захищеного повідомлення}%
  \label{fig:udp-encrypt-flowchart}
\end{figure}
